\noindent\textbf{Elección del procedimiento.} En fluido ilimitado el par por longitud es $4\pi\mu\Omega R^2$. Una carcasa fija multiplica ese par por $R_2^2/(R_2^2-R^2)$. La potencia se obtiene multiplicando cada par por la misma velocidad angular, manteniendo claro qué geometría corresponde a cada caso.

\begin{enumerate}

\item Calcular el par por longitud de la solución exterior.

\[\fitmath{\mathrm{Tfree}=4\cdot \pi\cdot \mu\cdot W\cdot {R}^{2}}\]

\[\fitmath{\mathrm{Tfree}=4\cdot \pi\cdot (0.08)\cdot (150)\cdot {(0.025)}^{2}\approx 0.09424778\,\mathrm{N}}\]

\item Aplicar el factor de confinamiento de cilindros coaxiales.

\[\fitmath{\mathrm{Twall}=\frac{\mathrm{Tfree}\cdot {R_{2}}^{2}}{{R_{2}}^{2}-{R}^{2}}}\]

\[\fitmath{\mathrm{Twall}=\frac{(0.09424777961)\cdot {(0.03)}^{2}}{{(0.03)}^{2}-{(0.025)}^{2}}\approx 0.3084473\,\mathrm{N}}\]

\item Potencia disipada por longitud sin carcasa.

\[\fitmath{\mathrm{Pfree}=\mathrm{Tfree}\cdot W}\]

\[\fitmath{\mathrm{Pfree}=(0.09424777961)\cdot (150)\approx 14.13717\,\mathrm{W/m}}\]

\item Potencia con la carcasa: no intercambiar resultados entre geometrías.

\[\fitmath{\mathrm{Pwall}=\mathrm{Twall}\cdot W}\]

\[\fitmath{\mathrm{Pwall}=(0.3084472787)\cdot (150)\approx 46.26709\,\mathrm{W/m}}\]

\end{enumerate}
